% Generated draft: requires mathematical and source review.
\documentclass[11pt]{article}
\usepackage[margin=25mm]{geometry}
\usepackage{fontspec}
\setmainfont{Latin Modern Roman}
\usepackage{amsmath,amssymb,mathtools,unicode-math}
\setmathfont{Latin Modern Math}
\usepackage{enumitem,needspace,xurl,hyperref}
\setlength{\parindent}{0pt}
\setlength{\parskip}{5pt}
\setlength{\emergencystretch}{3em}
\newcommand{\N}{\mathbb{N}}
\newcommand{\Z}{\mathbb{Z}}
\newcommand{\Q}{\mathbb{Q}}
\newcommand{\R}{\mathbb{R}}
\newcommand{\C}{\mathbb{C}}
\newcommand{\F}{\mathbb{F}}
\newcommand{\A}{\mathbb{A}}
\newcommand{\PP}{\mathbb{P}}
\newcommand{\E}{\mathbb{E}}
\newcommand{\eps}{\varepsilon}
\newcommand{\dd}{\,\mathrm{d}}
\begin{document}
\section*{UnsolvedMath ID: 11000017}
\textbf{Source problem code:} catalogue entry
\subsection{Definitions and mathematical statement}
Using the usual convention suggested by the automorphism-order discussion, let \(N(g)\) be the maximum order of the holomorphic automorphism group of a compact genus-\(g\) Riemann surface. The excerpt does not repeat the preceding definition of \(N\), so this convention requires confirmation. Set
\[H=\{g\in\Z:g\geq2,\ N(g)=8(g+1)\},\qquad
F_H(s)=\sum_{g\in H}g^{-s},\]
where \(g^{-s}=\exp(-s\log g)\) for \(s\in\C\). The principal problem is to determine the absolute-convergence abscissa
\[s_0=\inf\{\sigma\in\R:\sum_{g\in H}g^{-\sigma}<\infty\},\]
so that the series converges absolutely for \(\operatorname{Re}s>s_0\), and has a singularity at the real point \(s=s_0\). The singularity is not specified to be a pole, nor does the source require divergence of the series at the boundary point itself; a holomorphic-extension question is different from boundary convergence. The requested finite real boundary presupposes an appropriate infinite set \(H\); its infinitude/known context is not supplied in this excerpt. For a finite \(H\), the series would instead be entire and have abscissa \(-\infty\).

The following source discussion gives additional directions and qualifications. It reports Accola [Ac]: for \(3\mid g\),
\[N(g)\geq8(g+3),\]
with equality attained infinitely often, and asks for related bounds in other genera and the frequency of attainment. It also proposes imposing algebraic restrictions on automorphism groups, such as nilpotence, solvability or being a \(p\)-group, studying both upper and lower bounds and how often sharp bounds are attained. Its example is Zomorrodian [Zo]: if \(\operatorname{Aut}(X_g)\) is nilpotent, then
\[|\operatorname{Aut}(X_g)|\leq16(g-1),\]
and equality implies that \(g-1\) is a power of \(2\). The converse is not asserted. These variants are not silently identified with the specific set \(H\). The text ends by introducing a following question, which is absent.

The research summary mentions Larsen's frequency theorem for attainment of the Hurwitz bound \(84(g-1)\) as a model, without its formula or citation; that different attainment set does not specify \(F_H\).
\subsection{Short English statement}
Determine the Dirichlet-series convergence threshold for genera attaining the low-symmetry bound \(8(g+1)\), and retain the source's proposed frequency variants and restricted-automorphism-group bounds.
\subsection{Sources}
\begin{itemize}
\item Local UnsolvedMath registry, permanent ID \(11000017\), exact source code \texttt{catalogue entry}.
\item Farb, \emph{Problems on Mapping Class Groups and Related Topics}, Problem 2.15, PDF page 29: \url{https://www.math.uchicago.edu/~farb/papers/mcgbook.pdf}. Registry access 2026-07-29; pdf-font-structure extraction.
\item Accola [Ac] and Zomorrodian [Zo] are cited internally for the stated bounds; full entries are absent. The summary names Larsen for a Hurwitz-bound frequency theorem without a precise reference.
\end{itemize}
\end{document}
